\documentclass[11pt]{article}
\usepackage[margin=1in]{geometry}
\usepackage{amsmath,amssymb,amsthm,hyperref}
\newtheorem{theorem}{Theorem}
\newtheorem{lemma}{Lemma}
\newcommand{\CC}{\mathbb C}
\newcommand{\RR}{\mathbb R}
\newcommand{\diag}{\operatorname{diag}}
\newcommand{\MLD}{\operatorname{MLD}}
\newcommand{\eMLD}{\operatorname{eMLD}}
\title{The extremal maximum-likelihood degree bound\\
\large A solution of MathDB 373405}
\author{Principia Math In Progress}
\date{27 September 2026}
\begin{document}
\maketitle
\begin{abstract}
We prove that the extremal maximum-likelihood degree of a graph is at
most its Gaussian maximum-likelihood degree. The proof deforms each
generic complex extremal critical point into a nonsingular Gaussian
critical point. A second-order determinant expansion resolves the
degeneracy in the added diagonal parameters.
\end{abstract}

\section*{Task 1: statement and likelihood coordinates}
The target is Conjecture 4.12 of \cite{ACGR}, page 19.
Degrees count complex critical points for generic data, on the domain
where the relevant determinant is nonzero.
\begin{theorem}\label{main}
For every connected finite simple graph $G$,
\[
 \eMLD(G)\leq\MLD(G).
\]
The same bound holds for disconnected graphs under either the strict
completion definition in the source or its componentwise extension.
\end{theorem}

Let $G$ have vertices $1,\ldots,n$ and edge set $E$. We first take
$n\geq2$ and $G$ connected. Give its edges complex weights $w_e$ and set
\[
 L(w)=\sum_{e=\{i,j\}\in E}w_e(e_i-e_j)(e_i-e_j)^{\mathsf T},
 \quad A(w)=L(w)_{\widehat n,\widehat n},\quad \tau(w)=\det A(w).
\]
Here a hat means that the row and column are deleted. The map from
weights to $A$ is injective: a Laplacian with zero row sums is determined
by this principal submatrix, and its off-diagonal entries determine the
weights. On $\tau\ne0$, the extremal score equations are
\begin{equation}\label{ext}
 \partial_{w_e}\log\tau(w)=v_e\qquad(e\in E).
\end{equation}
Indeed, the pseudo-determinant of $L$ is $n\tau$, and for a hollow
variogram $\Gamma$,
\[
 \tfrac12\operatorname{tr}(L\Gamma)
   =-\sum_{\{i,j\}\in E}w_{ij}\Gamma_{ij}.
\]
Thus \eqref{ext} is exactly the critical system for the source's
surrogate likelihood, with $v_{ij}=\Gamma_{ij}$; see \cite[Section 4,
equation (16)]{ACGR}. It also describes the spanning-tree concentration
model in \cite{CMS}.

Every symmetric Gaussian concentration matrix with this graph's
off-diagonal zero pattern has unique coordinates
\[
 K=L(w)+\diag(s_1,\ldots,s_n),\qquad F(w,s)=\det K.
\]
For a symmetric sample matrix $S$, set
$u_i=S_{ii}$ and $v_{ij}=S_{ii}+S_{jj}-2S_{ij}$ on edges. These are
invertible linear coordinates on the specified sample entries, and
\[
 \operatorname{tr}(KS)=\sum_i s_i u_i+\sum_e w_e v_e.
\]
Consequently the Gaussian score equations are
\begin{equation}\label{gauss}
 \partial_{w_e}\log F=v_e,\qquad
 \partial_{s_i}\log F=u_i,\qquad F\ne0.
\end{equation}

For completeness, generic fibers of both score maps are finite and
nonsingular. At positive real edge weights, $A$ is positive definite.
The Hessian of $\log\det A$ in a nonzero weight direction $\dot A$ is
\[
 -\operatorname{tr}(A^{-1}\dot A A^{-1}\dot A)<0,
\]
because the coordinate map is injective. Thus the score map has a
nonzero Jacobian determinant. It is a dominant rational map between
spaces of the same dimension, generically finite and unramified over
$\CC$. The identical argument at a positive definite $K$ applies to
\eqref{gauss}. We may therefore choose $v$ such that \eqref{ext} has
exactly $r=\eMLD(G)$ distinct solutions, all with $\tau\ne0$ and
nonsingular Hessian $\nabla_w^2\log\tau$.

\section*{The determinant expansion}
Fix any $w$ with $\tau(w)\ne0$ and put $C=A^{-1}$. Extend $C$ to an
$n$-by-$n$ matrix $\widetilde C$ by giving it a zero last row and column.
Define
\begin{equation}\label{resistance}
 R_{ij}=\widetilde C_{ii}+\widetilde C_{jj}-2\widetilde C_{ij},
 \qquad Z=\sum_{i=1}^n z_i.
\end{equation}
All these functions are holomorphic near the chosen $w$.
\begin{lemma}\label{expansion}
As a polynomial identity in the parameter $t$, locally on $\tau\ne0$,
\begin{equation}\label{expand}
 F(w,tz)=t\tau Z+
 t^2\tau\sum_{i<j}R_{ij}z_i z_j+O(t^3).
\end{equation}
Here the remainder is divisible by $t^3$.
\end{lemma}
\begin{proof}
Write $\widehat z=(z_1,\ldots,z_{n-1})^{\mathsf T}$ and
$D=\diag(\widehat z)$. The zero row sums give
\[
 L=\begin{pmatrix}A&-A\mathbf1\\-\mathbf1^{\mathsf T}A&
                  \mathbf1^{\mathsf T}A\mathbf1\end{pmatrix}.
\]
The Schur complement of $A+tD$ in $L+t\diag(z)$ equals
\[
 tZ-t^2\widehat z^{\mathsf T}C\widehat z+O(t^3),
\]
while
\[
 \det(A+tD)=\tau\left(1+t\sum_{i<n}C_{ii}z_i+O(t^2)\right).
\]
Multiplying proves \eqref{expand}, since
\[
 Z\sum_{i<n}C_{ii}z_i-\widehat z^{\mathsf T}C\widehat z
     =\sum_{i<j}R_{ij}z_i z_j.
\]
\end{proof}

\section*{Lifting every extremal critical point}
Since $F(w,0)=0$, the quotient
\[
 P(t,w,z)=\frac{F(w,tz)}{t}
\]
is a polynomial, including at $t=0$. Consider the following $|E|+n$
equations in the same number of unknowns $(w,z)$:
\begin{align}
 H_e&=\partial_{w_e}P-v_eP=0 &&(e\in E),\label{He}\\
 H_0&=\partial_{z_n}P-P=0,\label{H0}\\
 H_i&=\frac{\partial_{z_i}P-\partial_{z_n}P}{t}=0
                         &&(1\leq i<n).\label{Hi}
\end{align}
The last expressions are polynomials because
$P(0,w,z)=\tau(w)Z$, whose $z_i$ derivatives all equal $\tau$.
When $t\ne0$ and $P\ne0$, these equations are equivalent to
\eqref{gauss} for
\begin{equation}\label{data}
 s_i=tz_i,\qquad u_i=1/t\quad(1\leq i\leq n),
 \qquad v_e\text{ as fixed above}.
\end{equation}

Let $w^*$ be any of the $r$ chosen extremal critical points. At $t=0$,
Lemma~\ref{expansion} turns \eqref{He}--\eqref{Hi} into
\begin{align*}
 H_e&=Z(\partial_{w_e}\tau-v_e\tau),\\
 H_0&=\tau(1-Z),\\
 H_i&=\tau\bigl((Rz)_i-(Rz)_n\bigr)
     =\tau\left(C_{ii}Z-2(C\widehat z)_i\right).
\end{align*}
Thus they have the unique solution in $z$ given by
\begin{equation}\label{zstar}
 Z=1,\qquad
 \widehat z^*=\tfrac12 A\,(C_{11},\ldots,C_{n-1,n-1})^{\mathsf T},
 \qquad z_n^*=1-\sum_{i<n}z_i^*.
\end{equation}

Use coordinates $(w,\widehat z,Z)$ and order the equations as
$(H_e)_{e\in E},(H_i)_{i<n},H_0$. Their Jacobian at
$(t,w,z)=(0,w^*,z^*)$ is block lower triangular. Its upper diagonal
block is
\[
 \tau\,\nabla_w^2\log\tau,
\]
and its lower diagonal block is
\begin{equation}\label{block}
 \begin{pmatrix}
 -2\tau C&\tau\,(C_{11},\ldots,C_{n-1,n-1})^{\mathsf T}\\
 0&-\tau
 \end{pmatrix}.
\end{equation}
Both are nonsingular. In the upper right block, derivatives with respect
to $z$ vanish because \eqref{ext} holds. This verifies the triangular
claim without any assumption on the signs of the complex weights or
on the entries of $z^*$.

The holomorphic implicit-function theorem now gives a solution
$(w(t),z(t))$ near each $(w^*,z^*)$ for all sufficiently small $t$.
Its Jacobian stays nonsingular, and
$P(t,w(t),z(t))\to\tau(w^*)\ne0$. Because the $r$ limiting points are
distinct, their neighborhoods may be chosen disjoint. For a common
sufficiently small nonzero $t$, we obtain $r$ distinct nonsingular
solutions of \eqref{gauss}. The changes of variables and equations
used above are invertible for this fixed $t\ne0$.

The special data in \eqref{data} need not themselves be generic.
Each of these $r$ nonsingular Gaussian solutions persists under a small
arbitrary perturbation of all data coordinates $(u,v)$, again by the
implicit-function theorem. The perturbation can be chosen in the
Zariski-open set where the number of solutions equals $\MLD(G)$.
It follows that $r\leq\MLD(G)$, proving the bound for connected
$G$ with $n\geq2$.

For a single vertex both degrees are one. If $G$ has two or more
connected components, every graph-supported Laplacian has at least two
independent kernel vectors, the component indicators. Thus it cannot
have rank $n-1$. The source's strict variogram model, and its regular
complex completion domain, are then empty, giving degree zero and a
trivial inequality. Alternatively, if disconnected extremal models are
defined componentwise, their degrees multiply, just as Gaussian degrees
do. Multiplication of the connected inequalities proves the bound for
that extension as well. This completes Theorem~\ref{main}.

\section*{Verification and proposed follow-up tasks}
The exact verifier checks the determinant coefficients and the two
Jacobian blocks on all connected labeled graphs with two through four
vertices, and additional graphs through seven vertices. These finite
calculations supplement the all-graph proof above.
\begin{description}
\item[Task 2 (proposed).] Classify connected graphs for which
$\eMLD(G)=\MLD(G)$, beyond the known chordal case.
\item[Task 3 (proposed).] Describe and count the Gaussian critical points
that do not converge to the extremal branches constructed here, in order
to express $\MLD(G)-\eMLD(G)$ geometrically.
\item[Task 4 (proposed).] Determine whether extremal ML degree is
monotone under connected induced subgraphs, as suggested by the
spanning-tree-model questions in \cite{CMS}.
\item[Task 5 (proposed).] Extend the deformation to concentration
spaces with a fixed kernel of dimension greater than one, identifying
conditions that make the analogue of \eqref{block} nonsingular.
\end{description}

\begin{thebibliography}{2}
\bibitem{ACGR} C. Am\'endola, J. I. Coons, A. Grosdos, and F. R\"ottger,
\emph{Algebraic statistics of H\"usler--Reiss graphical models in
multivariate extremes}, arXiv:2603.02191v1 (2026), especially
Conjecture 4.12. \url{https://arxiv.org/abs/2603.02191}.
\bibitem{CMS} S. Cox, P. Misra, and P. Semnani,
\emph{Homaloidal polynomials and Gaussian models of maximum likelihood
degree one}, Algebraic Statistics \textbf{15} (2024), 167--198.
\url{https://doi.org/10.2140/astat.2024.15.167}.
Final author version: \url{https://arxiv.org/abs/2402.06090v3}.
\end{thebibliography}
\noindent\textit{Review status.} Author-agent review; no independent peer
review or proof-assistant verification. The accompanying audits record
the source coverage and bounded prior-claim checks.
\end{document}
